\documentclass[10pt,a4paper]{amsart}
\usepackage[margin=19mm]{geometry}
\usepackage{amsmath,amssymb,mathtools}
\usepackage[scr=rsfs,cal=euler]{mathalfa}
\usepackage{xfrac}
\usepackage[unicode,hidelinks,bookmarks=false]{hyperref}
\newcommand{\ie}{i.e.\ }
\newcommand{\cf}{cf.\ }
\newcommand{\resp}{respectively\ }
\newcommand{\Subsection}{\subsection}
\providecommand{\nobreakdash}{\nobreak\hbox{-}}
\setlength{\emergencystretch}{2em}
\multlinegap0pt
\usepackage{amssymb}
\usepackage{enumerate}
\usepackage{mathtools}
\newtheorem{lem}{Lemma}
\newcommand{\ep}{\epsilon}
\title{Splitting singular fibers with periodic monodromies and their monodromy factorization}
\author{Hyunggi Kim}
\date{Source: arXiv:2608.24208v2 (CC BY 4.0)}
\begin{document}
\maketitle

\noindent\emph{Source and licence.} Splitting singular fibers with periodic monodromies and their monodromy factorization. Version arXiv:2608.24208v2 (CC BY 4.0), distributed by arXiv under the Creative Commons Attribution 4.0 International licence, http://creativecommons.org/licenses/by/4.0/, as recorded in the arXiv licence metadata for this exact version and shown on the abstract page https://arxiv.org/abs/2608.24208v2. Full licence notice: \texttt{licenses/arxiv-2608.24208-CC-BY-4.0.txt}.\par\smallskip

\noindent\textit{Editorial scope.} Counted unit: the article's lem lem:discriminant\_1 (source0.tex lines 1228--1233). The article's complete original proof of it is delivered with it (source0.tex lines 1235--1301, one \texttt{proof} environment of 67 lines). No mathematics is written, repaired or re-derived: the statement and the proof are the article's own lines, in the article's own notation, and no retained line is substituted or re-flowed.  No further source-local context is needed: every cross-reference of every retained line resolves inside the two delivered spans.  Nothing was omitted from any retained span, so the package carries no omission record.  No retained line contains a citation, so the wrapper carries no bibliography and no bibliography entry.  No retained line contains a figure environment, an image-inclusion command or drawing code, so no image is delivered and the package declares no asset dependency.  Editorial formatting only, in addition: an AMS \texttt{amsart} preamble that restates the article's own declarations and macros used by the retained lines, namely the environments \texttt{lem} on separate counters, together with the standard packages those lines need.  The retained environments print in this document as Lemma 1 in order of appearance, whereas the article prints the same items with the numbering of its own full document.  The complete native source of the article as distributed in the arXiv e-print for this exact version is bundled unchanged as \texttt{sources/208457/source0.tex}.
\begin{lem} \label{lem:discriminant_1}
The discriminant $\Delta_y(t)$ is given as
\begin{align}
\Delta_y(t)=(-1)^{\frac{n(n-1)}{2}} y^{2n^2-n}((2n)^{2n}{\ep_n^{2n(2n+1)}}-(2n+1)^{2n+1}y^n t^{2n}).
\end{align}
\end{lem}

\begin{proof}
Letting $a=-y^3-1$, we have $\psi_n(X,y)-t=(X+a)(y^n (X+a)^{2n}-{\ep_n^{2n}})-t$.
Since translation along the $X$-axis does not affect the discriminant, we can simply write 
\[F(X):=X(y^n X^{2n}-b)-t=y^n X^{2n+1}-bX-t=0,\]
where $b={\ep_n^{2n}}$.
Then $F'(X)=(2n+1)y^n X^{2n}-b$.
It is a well-known fact that the formula for the discriminant is given by 
\[\Delta_y(t)=(-1)^{\frac{n(n-1)}{2}} y^{-n} R(F,F'),\] 
where $R(F,F')$ is the resultant of $F$ and $F'$.
Recall that for polynomials $f(x)=\sum_{i=0}^n a_i x^i$ and $g(x)=\sum_{i=0}^m b_i x^i$, 
the resultant $R(f,g)$ is the determinant of the $(n+m)\times (n+m)$ Sylvester matrix:
\begin{align}
Syl(f,g)=
\begin{bmatrix} \label{eq:syl}
a_n & a_{n-1} &a_{n-2} & \cdots &a_{1}  & a_{0} & \cdots & 0 \\
0 & a_{n} &a_{n-1} & \cdots &a_{2}  & a_{1} & \cdots & 0 \\
\vdots & \vdots &\ddots & \vdots & \vdots &\vdots & \ddots &\vdots \\
0&0&\cdots&a_n&a_{n-1}&\cdots&\cdots&a_0 \\
b_m&b_{m-1}&\cdots&b_1&b_0&0&\cdots &0 \\
0&b_{m}&\cdots&b_2&b_1&b_0&\cdots &0 \\
\vdots &\vdots &\ddots&\vdots &\vdots &\vdots &\ddots &\vdots \\
0&0&\cdots&\cdots&b_m&b_{m-1}&\cdots&b_0
\end{bmatrix}
\end{align}
Substituting $f=F, g=F'$ into \eqref{eq:syl} yields the $(4n+1)\times (4n+1)$ matrix
\begin{align*}
Syl(F,F')=
\begin{bmatrix}
y^n & 0 & \cdots & \cdots & -b  & -t & \cdots & 0 \\
0 & y^n &0 & \cdots & 0  & -b & \cdots & 0 \\
\vdots & \vdots &\ddots & \vdots & \vdots &\vdots & \ddots &\vdots \\
0&0&\cdots&y^n&0&\cdots&\cdots&-t \\
(2n+1)y^n&0&\cdots&0&-b&0&\cdots &0 \\
0&(2n+1)y^n&\cdots&0&0&-b&\cdots &0 \\
\vdots &\vdots &\ddots&\vdots &\vdots &\vdots &\ddots &\vdots \\
0&0&\cdots&\cdots&(2n+1)y^n&0&\cdots&-b
\end{bmatrix}.
\end{align*}
Applying elementary row operations, we obtain
\begin{align*}
|Syl(F,F')|&=
\begin{vmatrix}
y^n & 0 & \cdots & \cdots & -b  & -t & \cdots & 0 \\
0 & y^n &0 & \cdots & 0  & -b & \cdots & 0 \\
\vdots & \vdots &\ddots & \vdots & \vdots &\vdots & \ddots &\vdots \\
0&0&\cdots&y^n&0&\cdots&\cdots&-t \\
0&0&\cdots&0&2nb&(2n+1)t&\cdots &0 \\
\vdots &\vdots &\ddots&\vdots &\vdots &\vdots &\ddots &\vdots \\
0&0&\cdots&\cdots&0&\cdots &2nb&(2n+1)t \\
0&0&\cdots&\cdots&(2n+1)y^n&0 &\cdots & -b
\end{vmatrix}\\
&=y^{2n^2} 
\begin{vmatrix}
2nb & (2n+1)t &0 & \cdots & 0  & 0  \\
0&2nb &(2n+1)t&\cdots&\cdots&0 \\
\vdots &\vdots &\ddots&\vdots &\vdots &\vdots \\
0&0&\cdots&\cdots &2nb&(2n+1)t \\
(2n+1)y^n&0&\cdots&\cdots &0& b
\end{vmatrix} \\
&=y^{2n^2}((2n)^{2n}b^{2n+1}-(2n+1)^{2n+1}y^n t^{2n}),
\end{align*}
so the discriminant $\Delta_y(t)$ is given by
\begin{align*}
\Delta_y(t)=(-1)^{\frac{n(n-1)}{2}} y^{2n^2-n}((2n)^{2n}({\ep_n^{2n}})^{2n+1}-(2n+1)^{2n+1}y^n t^{2n}),
\end{align*}
which proves the claim.
\end{proof}

\end{document}
